\documentclass[12pt]{article}
\usepackage[utf8]{inputenc}
\usepackage[T1]{fontenc}
\usepackage{hyperref}
\title{Teme Language Resources for K-12 Curriculum Development}
\author{Compiled by Claude Sonnet 5 (claude-sonnet-5) for Andr\'as Kornai}
\date{}
\begin{document}
\maketitle
\section{Introduction}
Teme (ISO 639-3 tdo; also Shoo, Tema, Temme) is an Adamawa language of the Yendang group, spoken in Adamawa State, north-eastern Nigeria (Mayo Belwa and Fufore LGAs). Wikipedia cites about 4,000 speakers (1995); Joshua Project gives 6,800--7,000 for the people group. Very little material exists.
\section{Primary Educational Resources}
No primers, textbooks or school materials were found in search results.
\section{Linguistic and Reference Materials}
\begin{itemize}
\item Wikipedia: \url{https://en.wikipedia.org/wiki/Teme_language}
\item Yendang languages (Wikipedia): \url{https://en.wikipedia.org/wiki/Yendang_languages}
\item Numerals of Teme (Max Planck EVA, Channumerals): \url{https://lingweb.eva.mpg.de/channumerals/Teme.htm}
\item ASJP word list entry: \url{https://asjp.clld.org/languages/TEME}
\end{itemize}
\section{Religious and Community Resources}
\begin{itemize}
\item Joshua Project, Teme people group (reports gospel recordings as the only Christian resources in Teme): \url{https://joshuaproject.net/people_groups/15332}
\item Scripture Earth page for tdo: \url{https://www.scriptureearth.org/00eng.php?iso=tdo}
\end{itemize}
\section{Digital and Audio Resources}
\begin{itemize}
\item Global Recordings Network, Teme: \url{https://globalrecordings.net/en/language/tdo}
\end{itemize}
\section{Curriculum Development Context}
No evidence of school use or an established orthography was found. A curriculum would begin with documentation and orthography work, using the word lists and recordings above.
\end{document}